\chapter{DOGMA checkpoints, branching and prefix state}

Three requests begin with the same prompt. Each has consumed tokens
$(1,2,1)$, but one will continue with token $0$ and another with token $2$.
Should the engine repeat the prompt computation? Should both requests point
to the same writable state? Neither choice is necessary. We can preserve one
immutable description of the prefix and restore independent mutable states.
The difficult part is specifying exactly what that description means.

Chapter 31 supplied an artifact-bound native recurrence; Chapter 32 supplied
owned storage with generational handles and in-flight pins. This chapter
connects them through a bounded, data-only checkpoint protocol. The executable
reference imports the actual Chapter 32 pool and Chapter 31 model. We derive
continuation equivalence, deliberately break it, then account for the storage
and replay costs of reusable prefixes.

\section{Cut the computation at a semantic boundary}
\index{Semantic boundary}

Let $x_{1:t}=(x_1,\ldots,x_t)$ be the ordered integer token prefix already
consumed, and $s_t$ the complete model state immediately after token $t$.
For fixed parameters $\theta$ and declared state semantics, write
\begin{align}
 s_{t+1}&=F_\theta(s_t,x_{t+1}),\\
 z_{t+1}&=G_\theta(s_{t+1}).
\end{align}
Here $z_{t+1}$ denotes logits after consuming the next token, not the token
sampled from those logits. A transport packet, scheduler iteration and
model-token boundary are different clocks. We snapshot only at the last one,
when no partially accepted native operation is outstanding.

\Plate{01-boundary}{A checkpoint cuts after the consumed prefix, not after a
network packet. The pending suffix is excluded. Position and any additional
state needed for continuation belong at the cut.}{fig:boundary}

For the restricted Chapter 31 reference, $s_t=(t,h_t)$ and
$h_t\in\mathbb{R}^D$. The model artifact fixes an embedding matrix
$E\in\mathbb{R}^{V\times D}$, coordinate retentions $a\in[0,1)^D$,
write coefficients $w\in\mathbb{R}^D$, and a head
$U\in\mathbb{R}^{D\times V}$. Its update is
\begin{equation}
 e=E_{x_{t+1},:},\quad
 h_{t+1}=a\odot h_t+\operatorname{sigmoid}(w\odot e)\odot e,
 \quad z_{t+1}=h_{t+1}^{\mathsf T}U.
 \label{eq:native}
\end{equation}
The printed vector orientation in the final product is notation; Python
represents $h$ as a one-dimensional tensor. Native arithmetic here uses
finite CPU float64 values. This is a concrete non-Transformer inference
component, not a trained complete DOGMA model.

\subsection{Tensor equality is not state sufficiency}

Consider a machine with hidden value $h=0$ and auxiliary phase
$p\in\{0,1\}$. On the next token it emits $h+(-1)^p$. Two snapshots containing
only $h$ are identical, yet one emits $1$ and the other $-1$. No hash of $h$
can recover the missing phase. A model with a finite convolution window,
unfinished memory block, positional clock or segment-dependent transition
requires those objects in its continuation state.

The reference does not serialize every future DOGMA architecture.
Its semantic identifier names the restricted causal diagonal state.
A new recurrence or additional state requires a new schema/semantic contract
and tested migration, not a permissive shape check. An offline
reverse-complement suffix state from Chapter 21 cannot silently become a causal
prefix state just because its tensor has the same dimensions.

\section{Derive the continuation obligation}
\index{Continuation equivalence}

Define $F_\theta^*(s,y)$ by repeated application of $F_\theta$ over a finite
suffix $y$. A snapshot decoder $D_C$ and encoder $E_C$ should satisfy
\begin{equation}
 D_C(E_C(s_t))\equiv s_t,
 \label{eq:roundtrip}
\end{equation}
where $\equiv$ means equality of all semantically relevant values under the
declared numerical interpretation. Request IDs and physical addresses are
excluded: restore must create new ownership.

\begin{proposition}[Continuation equivalence]
Suppose the transition/readout are deterministic, model artifact and
semantics are unchanged, and equation~\ref{eq:roundtrip} holds for complete
state. For any legal suffix $y_{1:k}$, original and restored execution produce
the same successive states and logits.
\end{proposition}
\begin{proof}
At suffix length zero the semantic states agree by
equation~\ref{eq:roundtrip}. Assume equality after $j$ suffix tokens.
Both executions next apply the same $F_\theta$ to equal states and equal
token $y_{j+1}$, so their next states agree. Applying the same $G_\theta$
gives equal logits. Induction covers all $0\le j\le k$.
\end{proof}

This proof locates the engineering obligations. It does not prove that
an arbitrary tensor dump contains complete state, that two floating-point
kernels use the same reduction order, or that sampled generation replays
without decoder state. Those are premises to check.
Here finite binary64 values round-trip, and restore uses the same native
operator as ordinary continuation. Independent rational tests check fixture
arithmetic. Another precision policy needs its own parity contract.

\subsection{If restore is approximate}
\index{Checkpoint approximation error}

Suppose a compressed checkpoint introduces hidden error
$\|\delta_t\|_\infty\le\epsilon$. In equation~\ref{eq:native}, the additive
write depends on input, not on $h_t$. Under exact arithmetic,
\begin{equation}
 \delta_{t+k}=a^{\odot k}\odot\delta_t,\qquad
 \|\delta_{t+k}\|_\infty\le\alpha^k\epsilon,\quad
 \alpha=\max_j a_j<1.
 \label{eq:error}
\end{equation}
If each subsequent update adds error at most $\eta$ in that norm, summing
the geometric series gives
\begin{equation}
 \|\delta_{t+k}\|_\infty
 \le\alpha^k\epsilon+\eta\frac{1-\alpha^k}{1-\alpha}.
\end{equation}
For $\alpha=1/2$, the limiting additive term is $2\eta$.
This bound is specific to the stated recurrence. State-dependent gates need
their own Jacobian/stability argument. Even a small logit perturbation can
change an argmax near a tie; this norm bound is not an identical-generation
guarantee.

\section{The envelope binds interpretation to bytes}
\index{Checkpoint envelope}\index{State semantic version}

Use the context tuple
\[
 \kappa=(\tau,m,\nu,\xi,D,V,\mathrm{dtype}),
\]
where $\tau$ is the tenant namespace, $m$ the model-content digest, $\nu$ the
state semantics and $\xi$ the tokenization contract identifier. $D$ and $V$
declare hidden width and vocabulary size. The body and envelope are
\begin{equation}
 B_t=(f,\kappa,t,x_{1:t},h_t),\quad
 C_t=(b_t,d_t),\quad b_t=\operatorname{encode}(B_t),\quad
 d_t=\operatorname{SHA256}(b_t).
 \label{eq:envelope}
\end{equation}
The format $f$ governs parsing; $\nu$ governs the meaning of parsed state.
Neither replaces the model artifact. SHA-256 supplies the digest, not a
producer-authentication protocol.\cite{nist}

\Plate{02-envelope}{A digest binds the envelope's bytes. Context and exact
tokens determine how the payload may be used. Trusted provenance remains
a separate boundary: a self-consistent envelope is not an authorized one.}{fig:envelope}

Chapter 31's artifact function hashes ordered vocabulary and parameter
content. The local $\xi$ hashes the ordered vocabulary mapping.
It is not a universal tokenizer: deployment must bind preprocessing,
normalization, special-token policy, adapters and any conditioning that
changes transitions. A display name such as \texttt{model-latest} is not
a content identity.

Tenant IDs here arrive through a trusted in-process caller. Including $\tau$
prevents accidental namespace reuse in the reference; it does not authenticate
an API request. Authorization must precede construction of the context.

\subsection{Three different meanings of validity}

\begin{enumerate}
\item \textbf{Byte integrity:} bytes match the supplied digest.
   Someone able to replace both can compute a new matching pair.
\item \textbf{Structural compatibility:} schema, dimensions, finite values,
   model, semantics, tenant and exact prefix satisfy the local contract.
\item \textbf{Replay consistency:} hidden values equal those obtained by
   executing the prefix from its prescribed initial state.
\end{enumerate}
The first two do not establish the third. Replace $h_1=0.625$ by $11$ and
recompute the digest: the object can remain well-formed. Optional
\texttt{replay\_audit}\index{Replay audit} recomputes the actual native prefix and detects this
fixture forgery, but costs the work caching aims to avoid. In ordinary use,
trusted production and storage supply the missing provenance. Neither replay
consistency nor hashing grants permission to read another tenant's state.

\section{Parse bounded data, not executable objects}

The decoder accepts immutable bytes and a lowercase hexadecimal digest.
It rejects bodies above 262,144 bytes before JSON parsing. The schema allows
at most 4,096 tokens and 256 hidden coordinates: teaching limits, not
measured production requirements.

Duplicate object names, unknown or missing fields, nonfinite values, booleans
where integers are required, out-of-vocabulary tokens, inconsistent position,
undeclared dtype and noncanonical encoding are rejected. Default Python JSON
accepts some forms this protocol refuses, including repeated names and
nonfinite constants.\cite{python-json} Integer hidden values must also be
exactly representable in binary64 before conversion. An overflow spelling such as
\texttt{1e999} also requires a finite-value check; rejecting the literal word
\texttt{Infinity} is insufficient.

The encoder sorts keys, removes optional whitespace, emits ASCII escapes and
forbids NaN. This is a Python wire profile, not universal cross-language
canonicalization. Interoperability needs an explicitly standardized encoding
and numerical representation.

\Listing{checked}

This printed function checks context and exact prefix after \texttt{decode}
has enforced schema and bytes. Whole functions are extracted from
\texttt{drafts/ch33/reference.py}; the parser and negative tests remain
available there even when not all code is printed.

Avoiding executable loading narrows one risk surface, but bounded JSON is
not a sandbox. Resource limits and trustworthy surrounding code still matter.
PyTorch's restricted-loading documentation similarly separates reduced
code-execution exposure from complete protection against denial of service
and downstream memory hazards.\cite{pytorch-load}
The reference calls neither pickle nor \texttt{torch.load}.

\section{Capture only a settled owned state}
\index{Quiescent snapshot}

Chapter 32's handle identifies pool instance, slot, allocation generation,
owner pair and artifact. Those physical handles do not belong in a portable
checkpoint. A restarted process might reuse an integer slot for an unrelated
request. Restoring the handle would confuse historical address with current
ownership.

Capture validates the actual pool handle, refuses pinned/retiring storage,
checks $t=|x_{1:t}|$, and copies immutable snapshot values into the envelope.
This conservative quiescence policy can later be replaced by a consistently
captured immutable version under an explicit synchronization protocol.

\Listing{capture}

Calls are serialized and exclusive. A private pool-validation helper checks
the pin ledger; this is a teaching bridge, not a public allocator API.
Generation/owner/artifact fences remain intact.
The supplied prefix is trusted bookkeeping: the pool stores position and
hidden values, not a transcript proving their origin. The replay falsifier
makes this limitation visible. Authoritative request records could strengthen
the interface without removing the need for trusted execution.

\section{Restore allocates; it does not resurrect}
\index{Fresh-owner restore}

Restoring to owner $o'$ creates a fresh handle $\ell'$:
\[
 \operatorname{read}(\ell',o')=(t,h_t),\qquad
 \operatorname{owner}(\ell')=o'.
\]
It may occupy another slot or a later generation of the old slot. The released
handle stays stale even when semantic values are reconstructed exactly.

For $K$ branches, validate the entire owner group/checkpoint and ensure $K$
free slots exist before allocating. Retiring pinned slots are not free.
Under serialized calls, ordinary validation and capacity failures leave the
pool unchanged.

\Listing{restore_many}

This is preflight failure atomicity, not a crash transaction. A process
failure midway through allocation or another thread's simultaneous allocation
requires rollback, reservation or different synchronization. The teaching
loop does not establish distributed or hardware atomicity.

\section{Work through two actual native branches}

Use Chapter 31's fixture with $a=(1/2,1/4)$ and $w=(0,0)$. Both gates are
$1/2$. Tokens $0,1,2$ have embeddings $(0,0),(1,0),(0,1)$ respectively.
Starting from zero:
\begin{align*}
 h_1&=(1/2,0) &&\text{after token }1,\\
 h_2&=(1/4,1/2) &&\text{after token }2,\\
 h_3&=(5/8,1/8) &&\text{after token }1.
\end{align*}
The left branch consumes $0$, giving $(5/16,1/32)$. The right consumes $2$,
giving $(5/16,17/32)$. Both positions become four; the checkpoint stays at three.

\Plate{03-branches}{Equal prefix values are copied into independently owned
slots. Suffix transitions change only their respective branches. The
checkpoint is an immutable value, not shared writable state.}{fig:branches}

\begin{center}\input{\Generated/branches.tex}\end{center}

These numbers come from the imported recurrence. Tests independently use
rational arithmetic, then enumerate every prefix of length zero through
three and every two-token suffix. For each case, replay from zero and
restore-plus-suffix must match. This is a finite implementation check,
distinct from the induction proof.

Fresh CPU tensors prevent writable aliases with source/sibling storage.
\index{Autograd!checkpoint boundary}
Loading scalar values also discards the old computation graph. Equal inference
logits do not establish resumed backpropagation. Training restart needs
parameters, optimizer, gradient-accumulation, RNG and data-order state beyond
this protocol.

\section{A prefix key is an index, not an acceptance decision}
\index{Prefix key}\index{Longest-prefix lookup}

Define
\[
 k_t=H\bigl(\operatorname{encode}(\kappa,x_{1:t})\bigr).
\]
Unlike the checkpoint digest, this key excludes $h_t$; it names the
prefix/context being sought. The cache can therefore detect conflicting
state bodies for one deterministic prefix instead of silently replacing them.

\Plate{04-prefix}{Longest-prefix matching uses exact token identity within
one context. Reuse stops at token three here; only the two remaining tokens
are executed. A shorter hit is valid but saves less work.}{fig:prefix}

For query $q_{1:T}$ choose
\[
 t^*=\max\{t\le T:\text{a validated checkpoint for }(\kappa,q_{1:t})\text{ exists}\}.
\]
Otherwise start from initial state. A zero-prefix checkpoint is a real hit,
so the API returns both length and an optional object; zero alone cannot
distinguish that hit from a miss.

\Needspace{16\baselineskip}
\Listing{PrefixCache.longest}

After a hash-table hit, the reference checks full envelope, context and token
identity. A test replaces the key function by a constant to force an indexing
collision: a wrong prefix is rejected. This tests the acceptance path, not
SHA-256's empirical collision resistance.
Sampling seeds are not required in this deterministic model-state key;
transition-changing conditioning is. Stochastic state updates need relevant
random state or an explicitly nondeterministic reuse contract.

\subsection{Isolation and bounded eviction}
\index{LRU eviction}\index{Tenant namespace}

The LRU budget sums immutable envelopes' wire lengths. Reuse updates recency.
An oversized object is rejected before eviction; otherwise oldest envelopes
are removed until it fits. This counts encoded bytes, not Python heap,
allocator capacity, device buffers or total resident memory.

Restored branches have independent pool values, so source-envelope eviction
cannot invalidate them. Physical sharing would instead require live branches
to retain shared pages. Cache and branch ownership are separate reachability
roots.

Different tenants receive different keys and cross-tenant restore is absent.
Timing leakage, malicious callers and authorization still need service-level
policy. Current vLLM documentation offers comparison: prefix-dependent block
hashes, extra identity inputs and optional salting for isolation.\cite{vllm-design}
This reference inherits no production guarantee from that precedent.

\section{Copy-on-write must refine the copy oracle}
\index{Copy-on-write refinement}

Suppose branches A and B initially map to immutable physical page $P$.
A write by A may allocate $P_A$, copy $P$, redirect A's mapping, then modify
$P_A$. B still observes original $P$.

\Plate{05-cow}{A possible copy-on-write representation beside the independent
copy oracle. Mapping publication, reference counts and page lifetime add
obligations. Physical COW is not implemented in this chapter.}{fig:cow}

Let $R$ expand physical representation into the logical independent states
visible through the API. For permitted operation $u$, require
\[
 R\bigl(\operatorname{COWstep}(P,u)\bigr)
 =\operatorname{Copystep}(R(P),u).
\]
Reads must agree. An A-only write changes only A's expanded state. After
release/eviction, remaining live mappings must resolve. Tests need partial
pages, two writers, copy-allocation failure and cancellation between copy
and mapping publication.

For single-page state size $S$ and $K$ branches, independent copies use $KS$
state bytes. If every branch writes that page, COW can also need roughly $KS$,
plus mappings. Smaller independently writable pages reduce copied bytes only
for suitable write patterns and add bookkeeping. Memory savings are
workload-dependent.

\section{Derive checkpoint spacing and the metadata bill}
\index{Checkpoint spacing}\index{Repeated-prefix metadata}

Let stream length be $T$ tokens, checkpoint interval $B$ tokens, and hidden
width $D$. With eight bytes per coordinate and
$n=\lfloor T/B\rfloor$ nonzero periodic checkpoints, state payload alone
uses $8Dn$ bytes.

Our envelope also stores each complete prefix. Under a separate logical
convention of four bytes per token ID, repeated prefix storage is
\begin{align}
 M_{\rm tokens}&=4\sum_{j=1}^{n}jB=2Bn(n+1)\ \text{bytes},\\
 M_{\rm payload}&=8Dn+2Bn(n+1).
 \label{eq:storage}
\end{align}
These are not actual JSON sizes: framing, decimal spellings, digests and
containers add bytes. The formula still exposes a common mistake:
fixed-width recurrent state does not make a cache of full-prefix envelopes
constant-memory. With $B=1$, copied token metadata grows quadratically in $T$.

If a desired boundary's residual offset from its preceding checkpoint is
uniform on $\{0,\ldots,B-1\}$, expected replay is
\begin{equation}
 \mathbb{E}[r]=\frac{1}{B}\sum_{r=0}^{B-1}r=\frac{B-1}{2}.
 \label{eq:replay}
\end{equation}
Uniformity is a workload assumption. Prefix junctions may cluster at templates
or block boundaries; finite-stream edges may change the residual distribution.

\Plate{06-cost}{Modeled storage and replay for $T=64$, $D=8$.
Token metadata counts repeated complete prefixes with four-byte IDs.
These curves are arithmetic, not measured latency or actual JSON bytes.}{fig:cost}

\begin{center}\input{\Generated/costs.tex}\end{center}

At $B=8$, eight hidden payloads use 512 bytes, copied prefixes 1,152 bytes,
and expected replay 3.5 tokens. At $B=1$, these totals rise to 4,096 and
8,320 bytes to eliminate replay. Chained hashes and a verified shared token
store can avoid transcript duplication, but introduce lifetime/integrity
dependencies.

The simple longest-prefix loop serializes successively shorter prefixes.
Since $\sum_{t=0}^{T}t=T(T+1)/2$, a miss can require $O(T^2)$ token-processing
work for keys despite expected constant-time dictionary lookup per candidate.
Tries or incremental indices can improve that cost without replacing
checkpoint validation.

\section{State snapshots, generation restart and durability}

A model-state checkpoint resumes $F_\theta$. Sampled-text restart may also
need RNG algorithm/state, sampler settings, grammar automaton, stop policy
and committed-output cursor. Training adds optimizer and data-order state.
These are different restart contracts; one tensor is not unqualified
resumability.

Our cache is process-local. Durable publication must specify survived failures.
One candidate sequence is: encode/validate; write an unadvertised object;
persist contents under the storage policy; publish its identity atomically;
make it discoverable. Checksums alone do not make writes atomic or durable.
Filesystem rename/directory persistence and object-store transactions require
platform-specific evidence; no such implementation is shipped here.

Current serving systems illustrate semantic granularity. vLLM documents
recurrent-state reuse at selected block boundaries. Additional checkpoints
at shared-prefix junctions require specific configurations.\cite{vllm-apc,vllm-args}
Those are comparison points, not evidence that this reference implements
those models or outperforms their engines.

\section{Exercises with worked reasoning}
\begin{enumerate}
\item \textbf{Predict a branch.} Continue the fixture with token $1$.
 \emph{Solution:} $(5/16+1/2,1/32)=(13/16,1/32)$ at position four.
\item \textbf{Prove a cut.} What is the induction's base case?
 \emph{Solution:} equality of complete semantic state, not physical address
 or only the hidden array.
\item \textbf{Omitted clock.} Construct equal arrays with unequal next outputs.
 \emph{Solution:} $h=0$, phase $p\in\{0,1\}$, output $(-1)^p$.
\item \textbf{Position.} Why keep $t$ if this recurrence does not use it
 algebraically? \emph{Solution:} to bind the consumed-token boundary and
 catch inconsistent request bookkeeping.
\item \textbf{Rehashed forgery.} Change $h$ and recompute SHA-256.
 \emph{Solution:} integrity/schema checks can pass; replay consistency and
 authorization remain separate.
\item \textbf{Boolean token.} Why not accept \texttt{True} as 1?
 \emph{Solution:} subtype convenience is not the declared exact-integer wire type.
\item \textbf{Overflow.} Is rejecting \texttt{NaN} sufficient?
 \emph{Solution:} no; \texttt{1e999} may parse to a nonfinite value.
\item \textbf{Fresh ownership.} Restore into a released slot.
 \emph{Solution:} the new generation does not reactivate the stale handle.
\item \textbf{Capacity.} Two slots exist; one is retiring with a pin.
 Can two branches restore? \emph{Solution:} no, until final acknowledgement.
\item \textbf{Group restore.} Why preflight all owners?
 \emph{Solution:} a later duplicate/unauthorized owner must not leave earlier
 branches allocated after ordinary validation failure.
\item \textbf{Autograd.} Does equal hidden state preserve gradients?
 \emph{Solution:} no; scalar serialization and fresh tensors cut the graph.
\item \textbf{Collision.} Force identical lookup keys for two prefixes.
 \emph{Solution:} full token/context comparison must reject the wrong envelope.
\item \textbf{Zero hit.} Why return an optional object as well as length?
 \emph{Solution:} zero can mean an initial-state checkpoint or no hit.
\item \textbf{Eviction.} Why can a restored branch outlive its source envelope?
 \emph{Solution:} it owns copied values. Shared pages would require references.
\item \textbf{Storage.} Verify $T=64,D=8,B=16$.
 \emph{Solution:} $n=4$; hidden 256 bytes, token IDs 640, total 896,
 expected replay 7.5 tokens under the declared model.
\item \textbf{Index cost.} Do $T$ dictionary probes imply $O(T)$ work?
 \emph{Solution:} not when encoding prefixes processes quadratically many IDs.
\item \textbf{Approximation.} Bound error for $\alpha=1/2$, $\epsilon=0.01$,
 $\eta=0$ after three steps. \emph{Solution:} at most $0.00125$ in hidden
 infinity norm; identical argmax is not established.
\item \textbf{COW falsifier.} Design a two-writer test.
 \emph{Rubric:} compare every read with independent copies; interleave distinct
 coordinate writes/releases; inject allocation failure; check mappings and
 live references during eviction.
\item \textbf{Benchmark.} Compare periodic and junction-aware checkpoints.
 \emph{Rubric:} fix model, precision, request traces and quotas; report reuse,
 replay, metadata-inclusive bytes, restore traffic and latency distributions.
 Include low-overlap and cache-thrashing controls.
\item \textbf{Restart.} Extend the protocol to sampled generation.
 \emph{Rubric:} enumerate RNG, sampler, constraint and output-commit state; specify
 failure boundaries, reject incompatible revisions, and test uninterrupted
 versus restarted traces. Durable publication requires crash tests.
\end{enumerate}

\Needspace{16\baselineskip}
\section{What the next engine transformation must preserve}

We have immutable inference snapshots, exact context/token binding,
fresh-owner restore, independent branches and bounded prefix caching.
Physical COW, authentication, durable publication and optimized device code
remain outside the executable copy oracle.

Chapter 34 will gather several owned states into a batched transition and
scatter results back. It must preserve each branch's values, exact token
position and ownership despite padding, cancellation and changing batch
membership. Throughput transformations are legal only when that per-request
contract survives.

\section*{Selective glossary}
\addcontentsline{toc}{section}{Selective glossary}
\textbf{Semantic boundary}: point after a complete declared transition.
\textbf{Checkpoint envelope}: bounded body plus integrity metadata.
\textbf{Continuation equivalence}: equal behavior on a common legal suffix.
\textbf{Prefix key}: context/token content index, not state acceptance.
\textbf{Copy-on-write refinement}: physical sharing that preserves independent-copy
observations.
\textbf{Replay audit}: recomputation consistency check, not authentication.
